% TOP_PROBLEM: {"id":30007698,"problem_number":"LOCAL-30007698","title":"Protection of AI outputs","category_id":21,"category":{"id":21,"name":"economics","display_name":"Economics","description":"Open economic theory statements, published research agendas, and proposed scoped specifications; question provenance and historical priority are qualified per record.","slug":"economics","order_index":21,"created_at":"2026-10-08T00:00:00Z"},"status":"provenance_pending","external_url":null,"legacy_ids":[],"published":false,"statement":"Compare sector-specific protection rules for AI-assisted works by discounted creation, access and cumulative-innovation welfare.","economics":{"original_id":187,"field":"Innovation and AI economics","kind":"design","formalization_basis":"proposed_specification","source_status":"not_verified","priority_status":"not_established","priority_note":"No exact open-problem passage was verified. The supporting publication does not establish priority or unresolved status.","earliest_verified_reference":null,"current_reference":{"authors":["Gaétan de Rassenfosse","Adam B. Jaffe","Joel Waldfogel"],"title":"Intellectual Property and Creative Machines","year":2024,"url":"https://www.nber.org/papers/w32698","doi":"10.3386/w32698","locator":"Primary abstract and indexed chapter conclusion snippets; complete agenda paragraph not retrieved","evidence_summary":"Addresses economic IP trade-offs for creative machines; exact compensation and output-protection agenda passages remain unverified."},"formalization_note":"The cited primary work supports the subject or supplies solutions in particular settings, but no exact open-question passage for this specification was verified. The mathematical target and restrictions are proposed here, with no canonical-conjecture or priority claim."},"statusQualification":"Provenance pending: an explicit published statement of this exact open question was not verified. This is a catalogue-authored candidate specification, not a certified open conjecture. The cited primary work supports the subject or supplies solutions in particular settings, but no exact open-question passage for this specification was verified. The mathematical target and restrictions are proposed here, with no canonical-conjecture or priority claim.","statusReviewedAt":"2026-10-08"}
% FIRST_POSTED: 2026-10-08
% ENTRY_KIND: statement_only
% !TeX program = lualatex
\documentclass[11pt]{article}
\usepackage{fontspec}
\usepackage[a4paper,margin=25mm]{geometry}
\usepackage{amsmath,amssymb,mathtools}
\usepackage{xurl}
\usepackage[hidelinks]{hyperref}
\newcommand{\sref}[2]{\hyperlink{source:#1:#2}{[S#2]}}
\setlength{\emergencystretch}{3em}
\begin{document}
% problemId:30007698
\section*{Protection of AI outputs}\label{30007698}
\subsection{Definitions and mathematical statement}
Consider a finite-horizon market for creative works, with periods \(t=0,\ldots,H\), discount factor \(\delta\in(0,1]\), human creators and firms using AI to produce works. A protection rule \(\ell\in\mathcal L\) specifies duration, permitted reuse and enforceable licensing rights for AI-assisted outputs; \(\mathcal L\) is a finite set of legally feasible policy alternatives fixed for the exercise. Creators and firms choose investments, output prices and licensing offers in a stated dynamic equilibrium. Let \(B_t(\ell)\) be aggregate consumption value, \(C_t(\ell)\) real creation and distribution cost, and \(K_t(\ell)\) enforcement cost. Define
\[
W(\ell)=E\!\left[\sum_{t=0}^{H}\delta^t\{B_t(\ell)-C_t(\ell)-K_t(\ell)\}\right].
\]
Licensing payments are transfers within this welfare calculation; prices influence access and investment through equilibrium choices. New creators may enter, and later works may use earlier works, so cumulative innovation must be represented rather than assumed independent across periods. Characterize or estimate which \(\ell\) maximizes \(W(\ell)\) for a specified creative sector, allowing uncertainty about demand, substitution and reuse technology. Distinguish increases in output volume from changes in originality and access. This is a scoped proposed policy-design problem supported by economic IP trade-offs, rather than a canonical unsolved theorem or a claim about what protection current law provides.

\par\textbf{Formulation and scope.} The cited primary work supports the subject or supplies solutions in particular settings, but no exact open-question passage for this specification was verified. The mathematical target and restrictions are proposed here, with no canonical-conjecture or priority claim.

\par\textbf{Open-question provenance.} An explicit published open-question statement was not verified. Sources below are supporting literature and must not be treated as a first listing.

\par\textbf{Historical priority.} No exact open-problem passage was verified. The supporting publication does not establish priority or unresolved status.

\subsection{Short English statement}
Compare sector-specific protection rules for AI-assisted works by discounted creation, access and cumulative-innovation welfare.

\subsection{Sources}
\begin{itemize}\raggedright
\item[\textnormal{[S1]}]\hypertarget{source:30007698:1}{} Gaétan de Rassenfosse, Adam B. Jaffe, Joel Waldfogel. 2024. Intellectual Property and Creative Machines. Primary abstract and indexed chapter conclusion snippets; complete agenda paragraph not retrieved.
\url{https://www.nber.org/papers/w32698}.
DOI: 10.3386/w32698.
{\small\textit{Supporting or current literature; see evidence qualification. Addresses economic IP trade-offs for creative machines; exact compensation and output-protection agenda passages remain unverified.}}
\end{itemize}
\end{document}
